\chapter{Project Requirements}
\label{ch:requirements}

This chapter describes the comprehensive functional and non-functional requirements for QCanvas and the QSim SDK, followed by use cases, external interface specifications, and system features that define the complete system behavior across all four development sprints.

\section{Functional Requirements}

\subsection{Conversion Module}
The conversion module requirements define how source circuits are parsed, normalized, validated, and translated into OpenQASM 3.0. They also capture the statistics and diagnostics required for educational feedback and regression analysis.
\begin{itemize}
    \item FR-C1: The system shall parse Cirq, Qiskit, and PennyLane circuits using AST-based parsing with runtime fallback.
    \item FR-C2: The system shall generate OpenQASM 3.0 code conforming to the supported subset per project scope \cite{cross2021openqasm}.
    \item FR-C3: The system shall handle measurement and reset instructions correctly.
    \item FR-C4: The system shall expose conversion statistics including qubits, depth, and gate counts.
    \item FR-C5: The system shall validate generated OpenQASM 3.0 against the supported subset and provide detailed diagnostics \cite{cross2017openqasm}.
    \item FR-C6: The system shall support gate mapping with parameter and broadcast handling.
\end{itemize}

\subsection{Simulation Module (QSim SDK)}
The simulation module requirements capture the execution behavior of the QSim backends. These requirements emphasize determinism, backend selection, and the metadata needed to interpret results. The QSim SDK is published on PyPI and operates independently of the QCanvas Web IDE.
\begin{itemize}
    \item FR-S1: The system shall support statevector, density matrix, and stabilizer simulation backends \cite{nielsen2010quantum}.
    \item FR-S2: The system shall allow configurable shots (1--10,000) and return counts and metadata.
    \item FR-S3: The system shall cache identical executions to improve performance.
    \item FR-S4: The system shall allow seed configuration for deterministic runs.
    \item FR-S5: The system shall provide backend capability discovery (max qubits, supported gates).
    \item FR-S6: The system shall return results with counts, execution time, and optional state snapshots.
    \item FR-S7: The QSim SDK shall be installable via \texttt{pip install qsim-sdk} and importable as a standalone Python library.
\end{itemize}

\subsection{Hybrid CPU--QPU Module}
The hybrid module requirements describe how QCanvas separates classical orchestration from quantum execution, keeping the frontend responsive while heavy computation is delegated to the QSim backend.
\begin{itemize}
    \item FR-H1: The system shall separate classical orchestration (CPU) from quantum circuit execution (QPU/Simulator) \cite{peruzzo2014variational}.
    \item FR-H2: The system shall implement data flow: Source $\rightarrow$ QASM $\rightarrow$ Execution $\rightarrow$ Results.
    \item FR-H3: The system shall provide WebSocket progress updates for long-running operations.
    \item FR-H4: The system shall enforce simulation-first execution with optional pluggable endpoints.
    \item FR-H5: The system shall provide resource-aware execution with limits and timeouts.
\end{itemize}

\subsection{Web IDE Module}
The Web IDE requirements focus on the user-facing editor, visualization, and project-management features that make the platform usable as an interactive learning environment.
\begin{itemize}
    \item FR-W1: The system shall provide code editing with syntax highlighting and diagnostics via Monaco Editor.
    \item FR-W2: The system shall display real-time D3.js circuit visualization and Chart.js results panels.
    \item FR-W3: The system shall support multiple example projects and templates.
    \item FR-W4: The system shall provide keyboard shortcuts for frequent actions (e.g., \texttt{Ctrl+Shift+R} to run).
    \item FR-W5: The system shall support multi-file projects with workspace management.
\end{itemize}

\subsection{Educational Intelligence Module}
\begin{itemize}
    \item FR-E1: The system shall provide hover-based contextual explanations for quantum gates, algorithms, and concepts within the Monaco editor.
    \item FR-E2: The system shall implement a gamification engine with XP rewards, level progression, and at least 44 unique achievements.
    \item FR-E3: The system shall allow users to publish, discover, and remix community circuit snippets with metadata (framework, difficulty, tags).
    \item FR-E4: The system shall track user learning progress with activity history and achievement dashboards.
\end{itemize}

\subsection{API and Realtime}
\begin{itemize}
    \item FR-A1: The system shall expose REST endpoints for conversion and validation.
    \item FR-A2: The system shall expose WebSocket endpoints for progress streaming.
    \item FR-A3: The system shall provide structured error payloads with codes and hints.
    \item FR-A4: The system shall provide discovery endpoints for frameworks and backends.
    \item FR-A5: The system shall support batch conversion and simulation operations.
\end{itemize}

\subsection{Validation}
\begin{itemize}
    \item FR-V1: The system shall validate OpenQASM 3.0 output against the allowed subset.
    \item FR-V2: The system shall suggest nearest supported constructs for incompatibilities.
    \item FR-V3: The system shall emit precise diagnostics with line/column, code, and remediation.
    \item FR-V4: The system shall check for excluded features (timing, pulse-level) with clear guidance.
\end{itemize}

\section{Non-Functional Requirements}

\subsection{Reliability}
\begin{itemize}
    \item \textbf{NFR-R1: Service Availability} \\
    The system containers (Backend, Database, Redis) shall report a ``healthy'' status (200 OK) on their health-check endpoints 99.9\% of the time during active operation.
    \item \textbf{NFR-R2: Audit Trail} \\
    The system shall record an audit log entry for 100\% of API requests resulting in 4xx or 5xx status codes.
\end{itemize}

\subsection{Usability}
\begin{itemize}
    \item \textbf{NFR-U1: Keyboard Efficiency} \\
    The system shall execute the current quantum circuit when the user triggers \texttt{Ctrl+Shift+R} (or \texttt{Cmd+Shift+R}).
    \item \textbf{NFR-U2: Visual Accessibility} \\
    The user interface shall utilize the Monaco Editor's high-contrast \texttt{vs-dark} theme standards to ensure text legibility.
    \item \textbf{NFR-U3: Hover Response} \\
    The ``Quantum Explain It'' hover provider shall respond within 200ms of cursor placement on a recognized symbol.
\end{itemize}

\subsection{Performance}
\begin{itemize}
    \item \textbf{NFR-P1: Conversion Latency} \\
    The API shall process circuit conversion requests (AST parsing to QASM generation) in under 1500ms for circuits with $<$100 gates.
    \item \textbf{NFR-P2: Simulation Throughput} \\
    The QSim backend shall complete execution of standard algorithms (e.g., Grover's Search n=2) in under 5 seconds.
    \item \textbf{NFR-P3: Visualization Render Time} \\
    The D3.js circuit visualization shall render updates in under 200ms (sub-200ms average observed in testing).
\end{itemize}

\subsection{Security}
\begin{itemize}
    \item \textbf{NFR-S1: API Rate Limiting} \\
    The API shall enforce a rate limit of \textbf{20 requests per minute} per client IP address for conversion endpoints.
    \item \textbf{NFR-S2: Data Integrity Protection} \\
    The system shall sanitize 100\% of user inputs via ORM parameterization to prevent SQL injection vulnerabilities.
    \item \textbf{NFR-S3: CIA Compliance} \\
    Sensitive API keys shall be encrypted at rest using AES-256 (Confidentiality), passwords hashed with Bcrypt and per-user salts (Integrity), and RBAC shall limit resource access by role (Availability).
\end{itemize}

\clearpage
\section{External Interface Requirements}

\subsection{User Interfaces}

\subsubsection{Web IDE Interface}
\begin{itemize}
    \item Monaco-based code editor with syntax highlighting for Qiskit, Cirq, PennyLane, and OpenQASM 3.0.
    \item Multi-panel layout: Editor, Circuit Visualization, Results, Console, Statistics.
    \item File tree navigation with create, rename, and delete operations.
    \item Theme switcher (dark/light modes) with responsive design for desktop and tablet.
\end{itemize}

\subsubsection{Results and Visualization}
\begin{itemize}
    \item Real-time D3.js 2D and 3D circuit diagram updated as the user types.
    \item Tabbed results display: Console, Output, Histogram, OpenQASM, Errors, Stats.
    \item Interactive measurement histograms using Chart.js.
    \item Export functionality for results (JSON, CSV) and QASM files.
\end{itemize}

\subsection{Software Interfaces}

\subsubsection{Quantum Framework APIs}
\begin{itemize}
    \item Qiskit API (version 0.44+): Circuit construction and manipulation \cite{qiskit}.
    \item Cirq API (version 1.0+): Circuit operations and gate definitions \cite{cirq_developers_2023}.
    \item PennyLane API (version 0.33+): Quantum node processing \cite{bergholm2018pennylane}.
\end{itemize}

\subsubsection{QSim SDK (PyPI)}
\begin{itemize}
    \item Installable via \texttt{pip install qsim-sdk}.
    \item Exposes \texttt{StatevectorBackend}, \texttt{DensityMatrixBackend}, and \texttt{StabilizerBackend} classes.
    \item Accepts OpenQASM 3.0 string input and returns structured result dictionaries.
\end{itemize}

\subsection{Communication Interfaces}

\subsubsection{REST API}
\begin{itemize}
    \item Protocol: HTTPS in production, HTTP in development.
    \item Format: JSON request/response bodies.
    \item Authentication: JWT tokens with RBAC.
    \item Base URL: \texttt{/api/v1/}
\end{itemize}

\subsubsection{WebSocket Interface}
\begin{itemize}
    \item Protocol: WSS in production, WS in development.
    \item Endpoint: \texttt{/ws}
    \item Use cases: Progress updates, real-time simulation notifications.
\end{itemize}

\clearpage
\section{Use Cases}

The following use cases describe the core interactions between users and the QCanvas system.

\subsection{Functional Requirements Traceability}

\begin{longtable}{p{0.5\textwidth} p{0.5\textwidth}}
\toprule
\textbf{FR Group} & \textbf{Primary Related Use-Cases} \\
\midrule
\endhead
Conversion (FR-C1--FR-C6) &
UC-02: Write Circuit Code \\
& UC-03: Convert to OpenQASM \\
& UC-06: Handle Validation Errors \\
& UC-07: Export OpenQASM \\
\addlinespace[0.5em]
Simulation (FR-S1--FR-S7) &
UC-04: Execute Circuit \\
& UC-05: Visualize Results \\
& UC-09: Query Backend Capabilities \\
& UC-10: Hybrid CPU--QPU Execution \\
\addlinespace[0.5em]
Hybrid CPU--QPU (FR-H1--FR-H5) &
UC-04: Execute Circuit \\
& UC-10: Hybrid CPU--QPU Execution \\
\addlinespace[0.5em]
Web IDE (FR-W1--FR-W5) &
UC-01: Launch Web IDE \\
& UC-02: Write Circuit Code \\
& UC-05: Visualize Results \\
& UC-08: Load Example Circuit \\
\addlinespace[0.5em]
Educational (FR-E1--FR-E4) &
UC-11: Hover Explanation \\
& UC-12: Earn Achievement \\
& UC-13: Share Circuit Recipe \\
\addlinespace[0.5em]
API and Realtime (FR-A1--FR-A5) &
UC-03: Convert to OpenQASM \\
& UC-04: Execute Circuit \\
& UC-09: Query Backend Capabilities \\
\addlinespace[0.5em]
Validation (FR-V1--FR-V4) &
UC-03: Convert to OpenQASM (error paths) \\
& UC-06: Handle Validation Errors \\
\bottomrule
\caption{Traceability between FRs and Use Cases}
\label{tab:fr-uc-traceability}
\end{longtable}

\subsection{UC-01: Launch Web IDE}
\begin{itemize}
    \item \textbf{Primary Actor:} User
    \item \textbf{Precondition:} Browser has network access to the QCanvas URL.
    \item \textbf{Main Flow:}
    \begin{enumerate}
        \item User navigates to the application URL.
        \item System loads the Monaco editor, sidebar, and results panel.
        \item System fetches the default ``Hello World'' circuit.
        \item System displays the ``Ready'' status indicator.
    \end{enumerate}
    \item \textbf{Postcondition:} IDE is initialized and waiting for input.
\end{itemize}

\subsection{UC-02: Write Circuit Code}
\begin{itemize}
    \item \textbf{Primary Actor:} User
    \item \textbf{Precondition:} UC-01 complete; Editor is active.
    \item \textbf{Main Flow:}
    \begin{enumerate}
        \item User selects a target framework (e.g., Qiskit) from the dropdown.
        \item User types quantum circuit code into the editor.
        \item System performs real-time syntax checking and updates D3.js circuit diagram.
    \end{enumerate}
    \item \textbf{Postcondition:} Valid or invalid source code exists in the editor buffer.
\end{itemize}

\subsection{UC-03: Convert to OpenQASM}
\begin{itemize}
    \item \textbf{Primary Actor:} User
    \item \textbf{Precondition:} Source code is present in the editor.
    \item \textbf{Main Flow:}
    \begin{enumerate}
        \item User clicks the ``Compile'' button.
        \item System sends source code to the Conversion API.
        \item API parses the code and generates OpenQASM 3.0 intermediate representation.
        \item System returns the QASM code and displays it in the output panel.
    \end{enumerate}
    \item \textbf{Alternate Flows:}
    \begin{enumerate}
        \item [3a] \textit{Syntax Error:} API returns line number; System highlights the error line in red.
        \item [3b] \textit{Unsupported Gate:} System returns gate-not-supported error with remediation hint.
    \end{enumerate}
    \item \textbf{Postcondition:} OpenQASM 3.0 code is available for simulation.
\end{itemize}

\subsection{UC-04: Execute Circuit}
\begin{itemize}
    \item \textbf{Primary Actor:} User
    \item \textbf{Precondition:} Valid OpenQASM 3.0 code exists.
    \item \textbf{Main Flow:}
    \begin{enumerate}
        \item User clicks ``Run'' and configures parameters (Shots=1024, Backend=Statevector).
        \item System submits job to the QSim backend.
        \item QSim executes the circuit and returns results.
        \item System pushes results to the frontend via WebSocket.
    \end{enumerate}
    \item \textbf{Postcondition:} Measurement counts and state data are available in the client.
\end{itemize}

\subsection{UC-05: Visualize Results}
\begin{itemize}
    \item \textbf{Primary Actor:} User
    \item \textbf{Precondition:} UC-04 completed successfully.
    \item \textbf{Main Flow:}
    \begin{enumerate}
        \item System receives JSON result payload.
        \item System renders a bar chart using Chart.js for measurement counts.
        \item System displays execution metadata (time, backend version, qubit count).
    \end{enumerate}
    \item \textbf{Postcondition:} Results are visible to the user.
\end{itemize}

\subsection{UC-06: Handle Validation Errors}
\begin{itemize}
    \item \textbf{Primary Actor:} User
    \item \textbf{Main Flow:}
    \begin{enumerate}
        \item System detects semantic error (e.g., unsupported gate).
        \item System squiggles the error line in red in the editor.
        \item System displays a ``Problems'' pane with error codes (e.g., ERR-020) and remediation suggestions.
    \end{enumerate}
    \item \textbf{Postcondition:} User is aware of the exact location and nature of the error.
\end{itemize}

\subsection{UC-11: Hover Explanation (Quantum Explain It)}
\begin{itemize}
    \item \textbf{Primary Actor:} User (learner)
    \item \textbf{Main Flow:}
    \begin{enumerate}
        \item User hovers cursor over a gate or algorithm keyword (e.g., \texttt{H}, \texttt{CNOT}, \texttt{VQE}).
        \item Monaco hover provider triggers \texttt{getHoverForSymbol} lookup.
        \item System renders a markdown popover with the gate matrix, Bloch-sphere implications, and a brief textbook explanation.
    \end{enumerate}
    \item \textbf{Postcondition:} User receives contextual quantum education without leaving the editor.
\end{itemize}

\subsection{UC-12: Earn Achievement}
\begin{itemize}
    \item \textbf{Primary Actor:} User
    \item \textbf{Main Flow:}
    \begin{enumerate}
        \item User completes a qualifying activity (e.g., 50 simulation runs, sharing a recipe).
        \item Gamification service evaluates the activity against 44 achievement criteria.
        \item System awards XP, updates the user's level, and displays an achievement notification.
    \end{enumerate}
    \item \textbf{Postcondition:} Achievement is recorded and displayed on the user's profile dashboard.
\end{itemize}

\subsection{UC-13: Share Circuit Recipe}
\begin{itemize}
    \item \textbf{Primary Actor:} User
    \item \textbf{Main Flow:}
    \begin{enumerate}
        \item User opens the \texttt{ShareModal} and annotates the snippet with title, difficulty, framework, and tags.
        \item System publishes the recipe to the community library.
        \item Other users can discover, fork, and run the shared circuit.
    \end{enumerate}
    \item \textbf{Postcondition:} Recipe is publicly visible in the community library.
\end{itemize}

\clearpage
\section{Use Cases Diagram}
This diagram outlines the primary actors and their interactions with QCanvas across core flows: conversion, simulation, educational features, and community sharing.
\begin{figure}[H]
    \centering
    \includegraphics[width=\linewidth]{FYP-UseCase.png}
    \caption{Use Case Diagram}
    \label{fig:usecase}
\end{figure}

\clearpage
\section{Domain Model Diagram}
This diagram shows the core entities (Circuit, ConversionRequest, SimulationJob, SimulationResult, User, Project, Achievement, Recipe) and key relationships that support conversion, simulation, and educational features.
\begin{figure}[H]
    \centering
    \includegraphics[width=\linewidth]{FYP-DomainModel.png}
    \caption{Domain Model}
    \label{fig:domain-model}
\end{figure}
